\section{Recoupling the complete physical source at fixed momentum angle}
\label{sec:v77-recoupling}
Use the unchanged $W=W^\varepsilon$ and one fixed exact label
$(m,R,n,k)$. Let $\Omega$ be its regular event, with the interior
critical points removed. The removed sets have zero original area.
Write $\mu_t$ for the coarea fiber of
$d\alpha\,dp/(4\pi c_*)$ on this event. Thus
$b^\varepsilon(t)=\int W\,d\mu_t$.

\begin{theorem}[A complete, cap-free two-observable disintegration]
\label{thm:v77-angular-source}
The physical measure $(L,\vartheta)_\#(W\one_\Omega\nu_R^*)$
has a nonnegative density $A^\varepsilon(t,\theta)$ relative to
$dt\,d\theta$ on $\R\times\mathbb S^1$. For almost every $(t,\theta)$,
\begin{equation}\label{eq:v77-angular-sum}
 A^\varepsilon(t,\theta)=
 \frac1{4\pi c_*}
 \sum_{\substack{x\in\Omega:\ L(x)=t,\ \vartheta(x)=\theta}}
                       \frac{W(x)}{|B_m(x)|\chi(x)}.
\end{equation}
The sum is over actual initial states, with every physical word and
the exact occupation, displacement, initial and terminal tests retained.
At fixed $m$ its fibers are finite; no estimate for their cardinality
as $m$ grows is used. One has
\begin{align}
 b^\varepsilon(t)&=\int_{\mathbb S^1}A^\varepsilon(t,\theta)\,d\theta,
                                      \label{eq:v77-angular-marginal}\\
 \int_\R\int_{\mathbb S^1}A^\varepsilon(t,\theta)\,d\theta\,dt
  &=\int_{\Pi_{n,k,m,R}}W\,d\nu_R^*\le1.
                                      \label{eq:v77-angular-total}
\end{align}
These equalities do not normalize the retained source a second time.
\end{theorem}
\begin{proof}
By \eqref{eq:v77-jacobian} the absolute Jacobian of the physical
map $(\alpha,p)\mapsto(L,\vartheta)$ is $|B_m|\chi$.
The area formula, on a finite circle atlas and a disjoint refinement
of regular inverse charts, proves \eqref{eq:v77-angular-sum}.
Overlaps are removed at the level of physical initial states. The
formula is consequently unchanged by relabeling a word chart or a
cap. Tonelli and ordinary disintegration prove the two identities.

For finiteness of the fibers at this fixed count, the explicit
finite physical graph retained in the supplement describes contacts,
first-hit tests, section decisions, and endpoint momenta by finitely
many semialgebraic conditions in rational angular charts. Add
$g=r(\cos\theta,\sin\theta)$, $r>0$, and $L=t$.
The nonzero Jacobian makes each such fiber discrete. A definable
discrete set is finite, uniformly in parameters for this fixed
formula. This assertion is a fixed-count fact; it is not a pressure
estimate for an increasing family of words.
\end{proof}

The inverse twist in this formula retains the optical bound
\begin{equation}\label{eq:v77-retained-optics}
 \frac1{|B_m|}\le\frac{R c_0c_m}{\tau_-}
       \frac{\sinh\eta}{\sinh(m\eta)}
       \prod_{i\in S}\frac{53c_i}{6+47c_i},\qquad
 \eta=\operatorname{arcosh}(53/47),
\end{equation}
for any selected internal set $S$. For clarity, its denominator is
obtained from $y_0=0$, $f_0=1$, $y_1=\ell_0$ and
$f_i=f_{i-1}+2y_i/(Rc_i)$, $y_{i+1}=y_i+\ell_i f_i$.
The corner Schur complement gives $|B_m|^{-1}=Rc_0c_m/y_m$.
Lowering lengths to $\tau_-$ and potentials to $2/R_+$ gives the
hyperbolic-sine solution. Restoring selected potentials one at a time
gives the factors in \eqref{eq:v77-retained-optics} by their positive
scalar Schur complements. The complete proof is retained in the
supplement. We use its per-word estimate, not an unproved exponential
word-count bound. Identity \eqref{eq:v77-angular-total} is an actual
complete weighted summability identity, but only in integrated form.

Disintegrate the weighted level source also over $\theta$:
\begin{equation}\label{eq:v77-angular-conditionals}
              W\mu_t=\int_{\mathbb S^1}
               A^\varepsilon(t,\theta)\,\rho_{t,\theta}\,d\theta,
\end{equation}
where $\rho_{t,\theta}$ is a probability on physical initial states
with this roof, angle, and exact label whenever its coefficient is
positive. In the finite regular fiber it assigns weights proportional
to the summands of \eqref{eq:v77-angular-sum}. Its value where the
coefficient is zero is immaterial.

\begin{theorem}[Maximal recoupling with the physical angle retained]
\label{thm:v77-recoupling}
For almost every $(t,h)$ with $h>0$, the measure
\begin{equation}\label{eq:v77-recoupling}
 \Gamma_{t,h}=\int_{\mathbb S^1}
       \min\{A^\varepsilon(t,\theta),A^\varepsilon(t+h,\theta)\}
          (\rho_{t,\theta}\otimes\rho_{t+h,\theta})\,d\theta
\end{equation}
is a subcoupling of the complete physical weighted fibers at $t$ and
$t+h$. Among all subcouplings which preserve $\vartheta$, its matched
mass is maximal. Its unmatched initial mass is exactly
\begin{equation}\label{eq:v77-angular-defect}
 \Delta_h^\varepsilon(t)=\int_{\mathbb S^1}
          [A^\varepsilon(t,\theta)-A^\varepsilon(t+h,\theta)]_+\,d\theta.
\end{equation}
It follows that
\begin{equation}\label{eq:v77-angular-drop}
 [b^\varepsilon(t)-b^\varepsilon(t+h)]_+
                   \le\Delta_h^\varepsilon(t)\le R^+_{X,h}(t),
\end{equation}
where $R^+_{X,h}$ is the unmatched mass of the stopped physical-flow
subcoupling for the explicit field in Theorem~\ref{thm:v77-field}.
\end{theorem}
\begin{proof}
Each marginal of \eqref{eq:v77-recoupling} is dominated by the
corresponding measure in \eqref{eq:v77-angular-conditionals}.
For any other angle-preserving subcoupling, push its matched measure
to the common angle. Its density is bounded above by both
$A^\varepsilon(t,\theta)$ and $A^\varepsilon(t+h,\theta)$, and
therefore by their minimum. Formula \eqref{eq:v77-recoupling}
attains this upper bound. Subtracting its mass from the source mass
gives \eqref{eq:v77-angular-defect}, and target domination proves
the first inequality in \eqref{eq:v77-angular-drop}.

The stopped-flow graph for $X$ preserves $\vartheta$ by
\eqref{eq:v77-field-defining}. On its surviving part its source
weight is the minimum of the source density and the transported
target density, with the physical coarea Jacobian. The usual flow
change of variables gives domination of both physical marginals.
It is therefore one admissible competitor in the preceding maximality
argument. Its unmatched mass is $R^+_{X,h}$, which proves the last
inequality. Joint disintegration supplies almost-everywhere versions
in $(t,h)$; no pointwise selection on an exceptional roof is required.
\end{proof}

The optimization in this theorem is not over unrestricted scalar
couplings. It retains the explicitly specified physical angle. It
may pair two different regular words, but only when both are actual
states of the same original exact event at their respective roofs.
It does not splice a nonphysical orbit or pool different labels.
The full sum in \eqref{eq:v77-angular-sum} is formed before its
positive difference is taken. Thus a true word exit may be balanced
by actual incoming mass at the same angle and label; any remaining
loss is still charged. Such balancing was unavailable to a stopped
single-orbit graph. For example, two nonnegative angular profiles
supported on successive roof intervals can have a continuous sum
although the first profile has an exit. This is an algebraic
illustration of the distinction, not a claim about a particular
billiard trajectory. No comparison with every earlier scalar
short-excursion modulus is asserted.

\begin{proposition}[The angular residue at an interior critical roof]
\label{prop:v77-angular-residue}
Let $x_*$ be a regular interior critical point with a neighborhood
in the exact event, $F_*=F(x_*)$, and $K_*=D^2F(x_*)$.
Write $a_*=W(x_*)/(4\pi c_*|B_m(x_*)|)$. Its local angular
source has right-minus-left jump
\begin{equation}\label{eq:v77-angular-residue}
 j_*(\theta)=\frac{a_*}
       {e_\theta^{\mathsf T}\operatorname{cof}(K_*)e_\theta}\ge0,
 \qquad
 \int_{\mathbb S^1}j_*(\theta)\,d\theta
     =\frac{W(x_*)}{2c_*|B_m(x_*)|\sqrt{\det K_*}}.
\end{equation}
Every direction is an entrance, not a forward exit. This is a local
charge statement; no sign is assigned to other physical boundaries
at the same roof.
\end{proposition}
\begin{proof}
Translate $q_*$ to zero. The inverse function theorem for
$g=\nabla F$ gives a smooth inverse $q(r e_\theta)$ for small $r$.
Along it $dF/dr=r e_\theta^{\mathsf T}K^{-1}e_\theta>0$.
Thus, for each angle, the local branch starts at $F_*$ from above.
The density in \eqref{eq:v77-angular-sum} tends to
$a_*/(e_\theta^{\mathsf T}\operatorname{cof}(K_*)e_\theta)$.
The convergence is uniform in angle on this fixed neighborhood.
Diagonalizing the positive matrix $K_*$ and integrating
$(a\cos^2\theta+b\sin^2\theta)^{-1}$ gives
$2\pi/\sqrt{ab}$, which proves the formula. The neighborhood
is used only to calculate the residue. The globally defined field
and its uniform motion bounds use no such radius.
\end{proof}
